\documentclass[10pt]{article}
\usepackage[margin=0.8in]{geometry}
\usepackage{amsmath}
\usepackage[most]{tcolorbox}
\usepackage{listings}
\usepackage{xcolor}

\parindent=0pt

\definecolor{codebg}{RGB}{245,245,240}      % code background
\definecolor{filebg}{RGB}{220,220,215}      % filename bar
\definecolor{framecol}{RGB}{170,170,170}    % border color
\definecolor{commentcol}{RGB}{100,100,100}      % comment color
\lstdefinestyle{mycode}{
    language=C,
    basicstyle=\ttfamily\small,
    commentstyle=\color{commentcol},
    columns=fullflexible,
    keepspaces=true,
    showstringspaces=false,
    breaklines=true,
    tabsize=4
}
\newtcblisting{codefile}[2][]{%
    enhanced,
    colback=codebg,
    colframe=framecol,
    boxrule=0.5pt,
    arc=2pt,
    left=6pt,
    right=6pt,
    top=6pt,
    bottom=6pt,
    listing only,
    listing options={style=mycode},
    title=\texttt{#2},
    colbacktitle=filebg,
    coltitle=black,
    fonttitle=\small\ttfamily,
    attach boxed title to top left={
        yshift=-2mm,
        xshift=0mm
    },
    boxed title style={
        boxrule=0.5pt,
        colframe=framecol,
        colback=filebg,
        arc=2pt
    },
    #1
}


\title{CSE 4300: Prog 4}
\author{Gordon Chen, gbc24001}
\date{}

\begin{document}
\maketitle


\begin{codefile}{synch.h}
struct lock {
    char *name;
    volatile int locked;
    volatile struct thread *owner;
};
\end{codefile}
The \verb|lock| struct needs a boolean \verb|locked| flag and a pointer to the current lock \verb|owner|.\\


\begin{codefile}{synch.c}
struct lock *
lock_create(const char *name)
{
    struct lock *lock;

    lock = kmalloc(sizeof(struct lock));
    if (lock == NULL) {
        return NULL;
    }

    lock->name = kstrdup(name);
    if (lock->name == NULL) {
        kfree(lock);
        return NULL;
    }

    lock->locked = 0;
    lock->owner = NULL;
    return lock;
}
\end{codefile}
Initialize the lock to unlocked and no owner.\\


\begin{codefile}{synch.c}
void
lock_destroy(struct lock *lock)
{
    assert(lock != NULL);

    int spl = splhigh();
    assert(lock->locked == 0);
    assert(lock->owner == NULL);
    assert(thread_hassleepers(lock) == 0);
    splx(spl);

    kfree(lock->name);
    kfree(lock);
}
\end{codefile}
Check that the lock is unlocked and that no threads are sleeping on the lock before destroying.\\


\begin{codefile}{synch.c}
void
lock_acquire(struct lock *lock)
{
    assert(lock != NULL);
    assert(in_interrupt==0);

    int spl = splhigh();
    while (lock->locked == 1) {
        thread_sleep(lock);
    }
    assert(lock->locked == 0);
    assert(lock->owner == NULL);
    lock->locked = 1;
    lock->owner = curthread;
    splx(spl);
}
\end{codefile}
Sleep until the lock is unlocked. Then lock and set the owner to the current thread.\\


\begin{codefile}{synch.c}
void
lock_release(struct lock *lock)
{
    assert(lock != NULL);

    int spl = splhigh();
    assert(lock->locked == 1);
    assert(lock->owner == curthread);
    lock->locked = 0;
    lock->owner = NULL;
    thread_wakeup(lock);
    splx(spl);
}
\end{codefile}
Check that the thread trying to release is the thread that has the lock.
Then unlock and wakup all threads that are sleeping on the lock.\\


\begin{codefile}{synch.c}
int
lock_do_i_hold(struct lock *lock)
{
    assert(lock != NULL);

    int spl = splhigh();
    int ihold = (lock->locked) && (lock->owner == curthread);
    splx(spl);

    return ihold;
}
\end{codefile}
The current thread holds the lock if the lock is locked and the owner is the current thread.\\


\begin{codefile}{synch.c}
struct cv *
cv_create(const char *name)
{
    struct cv *cv;

    cv = kmalloc(sizeof(struct cv));
    if (cv == NULL) {
        return NULL;
    }

    cv->name = kstrdup(name);
    if (cv->name==NULL) {
        kfree(cv);
        return NULL;
    }
    return cv;
}
\end{codefile}
I left \verb|cv_create| and \verb|struct cv| unchanged. I don't think any more fields in \verb|cv| are
required.\\


\begin{codefile}{synch.c}
void
cv_destroy(struct cv *cv)
{
    assert(cv != NULL);

    int spl = splhigh();
    assert(thread_hassleepers(cv) == 0);
    splx(spl);

    kfree(cv->name);
    kfree(cv);
}
\end{codefile}
Check that no threads are sleeping on the cv before destroying.\\


\begin{codefile}{synch.c}
void
cv_wait(struct cv *cv, struct lock *lock)
{
    assert(cv != NULL);
    assert(lock != NULL);
    assert(in_interrupt==0);

    int spl = splhigh();
    assert(lock->locked == 1);
    assert(lock->owner == curthread);
    lock_release(lock);
    thread_sleep(cv);
    lock_acquire(lock);
    splx(spl);
}
\end{codefile}
Check that the current thread has the lock, then release the lock, and sleep on the cv.
When the thread wakes up, acquire the lock.\\


\begin{codefile}{synch.c}
void
cv_signal(struct cv *cv, struct lock *lock)
{
    assert(cv != NULL);
    assert(lock != NULL);

    int spl = splhigh();
    assert(lock->locked == 1);
    assert(lock->owner == curthread);
    thread_wakeup_one(cv);
    splx(spl);
}
\end{codefile}
Check that the current thread has the lock. Then wake up 1 thread sleeping on the cv.
\verb|thread_wakeup_one| is a new function I defined in \verb|thread.c|, see below.\\


\begin{codefile}{synch.c}
void
cv_broadcast(struct cv *cv, struct lock *lock)
{
    assert(cv != NULL);
    assert(lock != NULL);

    int spl = splhigh();
    assert(lock->locked == 1);
    assert(lock->owner == curthread);
    thread_wakeup(cv);
    splx(spl);
}
\end{codefile}
Check that the current thread has the lock. Then wakeup all threads sleeping on the cv.\\


\begin{codefile}{thread.h}
/*
 * Cause at most 1 thread sleeping on the specified address to wake up.
 * Interrupts must be disabled.
 */
void thread_wakeup_one(const void *addr);
\end{codefile}

\begin{codefile}{thread.c}
/*
 * Wake up at most one thread sleeping on "sleep address"
 * ADDR.
 */
void
thread_wakeup_one(const void *addr)
{
    int i, result;

    // meant to be called with interrupts off
    assert(curspl>0);

    // This is inefficient. Feel free to improve it.

    for (i=0; i<array_getnum(sleepers); i++) {
        struct thread *t = array_getguy(sleepers, i);
        if (t->t_sleepaddr == addr) {

            // Remove from list
            array_remove(sleepers, i);

            /*
             * Because we preallocate during thread_fork,
             * this should never fail.
             */
            result = make_runnable(t);
            assert(result==0);
            return;
        }
    }
}
\end{codefile}
Wake up a single thread for \verb|cv_signal|. This same as \verb|thread_wakeup| except that instead of
waking up all sleeping threads, it returns after waking up 1 thread.


\end{document}
